\section{实验设置}
\subsection{数据划分与预处理}
CIFAR-10有50000张训练图和10000张测试图，每张为32$\times$32的RGB图像，共10类\cite{krizhevskyCifar}。基础实验沿用A1分层划分：45000张训练、5000张验证，种子20261011，每类分别4500与500张。五折划分种子20261102，每折每类1000张，单折训练40000张、验证10000张。最终选定配置另在全部50000张训练图上重训。

像素先除255，再按通道标准化。通道均值与标准差分别为
\[ \boldsymbol\mu=(0.4914,0.4822,0.4465),\qquad \boldsymbol\sigma=(0.2470,0.2435,0.2616). \]
标准化常数取自完整训练池。增强时补4像素边界，随机裁剪回32$\times$32并以0.5概率水平翻转；填充值为标准化后的零。Cutout边长为16，中心随机。

\subsection{基准网络与优化参数}
基准结构见图\ref{fig:architecture}。三个卷积段各含两层3$\times$3卷积，通道为64/128/256；三次池化将空间尺寸降至4$\times$4，分类器为4096$\to$256$\to$10。无归一化模型有2196810个参数；归一化版本移除卷积bias并加入affine，共2197706个参数。
\begin{figure}[htbp]\centering
\includegraphics[width=\linewidth]{architecture.pdf}
\caption{基准CNN的张量变化与模块位置。叠片示意通道组，尺寸为通道$\times$高$\times$宽。}\label{fig:architecture}
\end{figure}

默认训练30轮，batch=128，峰值学习率0.02，采用SGD、动量0.9、Nesterov及weight decay=$5\times10^{-4}$。前2轮线性预热，之后余弦退火至0。基础实验使用原图和硬标签，增强与标签平滑在B1中逐项加入。

\subsection{实验矩阵与超参数选择}
\begin{table}[htbp]\centering\small
\caption{正式实验设计。次数为独立训练任务数；“三种子”指0、1、2，其余为种子0。}\label{tab:design}
\begin{tabularx}{\linewidth}{lXrr}
\toprule 组 & 内容 & 次数 & 轮数\\\midrule
E0 & 无归一化、无Dropout基准，三种子 & 3 & 30\\
E1 & FC概率0.3/0.5/0.7；卷积0.1/0.2与FC0.5组合，三种子 & 15 & 30\\
E2归一化 & BN/LN/IN/GN，三种子 & 12 & 30\\
E2批大小 & BN/GN/LN，batch=16且学习率降至0.0025 & 3 & 30\\
E2学习率 & none/BN/GN与0.005/0.1/0.4组合 & 9 & 30\\
E2组合 & BN+FC0.5，BN+卷积0.1+FC0.5，三种子 & 6 & 30\\
E3搜索 & 24组随机超参数 & 24 & 15\\
E3交叉验证 & 搜索前4名，各做5折 & 20 & 30\\
E3最终 & CV均值最优配置，全50000张重训 & 1 & 30\\
B1 & 增强、Cutout、标签平滑、加宽、架构替换 & 6 & 40\\
B2深度 & Plain/ResNet各20/32/56层 & 6 & 30\\
B2平移 & MaxBlurPool，无增强/裁剪翻转 & 2 & 40\\\bottomrule
\end{tabularx}\end{table}

随机搜索种子20261103，从$[0.005,0.2]$和$[10^{-5},2\times10^{-3}]$按对数均匀抽取学习率与weight decay，从$\{0,0.3,0.5\}$、$\{64,128,256\}$和$\{\mathrm{BN},\mathrm{GN}\}$抽取FC Dropout、batch和归一化。24组各训练15轮，按末轮验证准确率选前4名；随后从头训练每个候选的五折模型，以折均值选配置，最后在全部训练池重训30轮。

B1在CV最优配置上做40轮累计配方比较，依次加入裁剪翻转、Cutout、标签平滑和VGG加宽（128/256/512通道），最后一项为ResNet-56架构替换，采用16/32/64通道及全局平均池化分类头。B2深度实验统一BN、学习率0.1与裁剪翻转，比较同主路径深度有无捷径。两类集成均为softmax概率平均：一类由CV赢家的五个折模型组成，另一类从B1按验证准确率选3个成员。

\subsection{评价指标与模型选择}\label{sec:boundaries}
Accuracy衡量分类正确率，NLL衡量真实类别的预测概率，ECE衡量置信度与正确率的偏差。Accuracy与ECE用百分比表示，NLL取自然对数。三种子实验报告均值$\pm$样本标准差（ddof=1），五折验证报告折均值与折间标准差，其余为单种子结果。

随机搜索按验证准确率选候选，五折均值用于内部选型；集成成员和温度参数也由验证数据确定。全部正式训练完成后，统一进行完整测试集评估。
